% !Mode:: "TeX:UTF-8"
\fancyhead[CO]{\song\wuhao ~\chaptername \quad Transformer 模型关键算子的代码生成~}
\fancypagestyle{plain}{\fancyhead[CO]{\song\wuhao ~\chaptername \quad Transformer 模型关键算子的代码生成~}}

\chapter{Transformer 模型关键算子的代码生成}

% ########################################################################################
%                                        1 引言
% ########################################################################################
\section{引言}
伴随 Transformer 模型在大规模预训练、推理场景中获得广泛应用，多头注意力机制和矩阵乘法等关键算子的性能已成为影响整体模型执行效率的重要因素。尽管现有深度学习框架主要依靠厂商提供的高性能库以获得较好的算子实现，但这种方式在跨平台适配、结构灵活性与可扩展性等方面仍然面临较多挑战。特别是在异构硬件环境不断增多的情况下，各类 GPU 和国产加速芯片在执行模型、线程组织方式、存储层级结构及专用指令支持等方面存在差异，使得单一平台上的最优实现常常难以直接迁移到其他平台。

从编译技术角度看，当前主流深度学习编译框架虽然能够完成算子级别的自动生成和优化，但是仍然存在以下问题。首先，不少中间表示系统难以显式表达块级并行结构和数据分块策略，使得算子调度逻辑和硬件资源映射过程分散在不同优化阶段，无法进行统一建模和整体分析。其次，在矩阵乘和注意力等计算密集型算子中，性能瓶颈主要来自数据在多级存储结构之间的搬运和复用效率不足，算术运算本身并非主要限制因素；然而传统中间表示在描述共享内存布局、线程协作方法和访存重排策略等方面还不够理想，使复杂的访存优化难以在统一模式下完成。最后，在算子优化过程中，主要参数共同组成高维调优空间，如果缺少结构化参数策略和整体剪枝方式，自动调优会遇到组合爆炸问题，难以在合理时间内获得高性能实现。

\begin{figure}[htbp]
\centering
\includegraphics[width=0.95\textwidth]{3_1}
\caption{DeepGen 总体架构}\label{3_1}
\vspace{\baselineskip}
\end{figure}

本章针对上述问题，提出并实现了一种基于 MLIR 的 Transformer 主要算子代码生成与优化方法。本文通过面向 GPU 块级执行结构的分层 IR 表达形式，对算子的结构、访存路径、并行映射进行统一建模。图 \ref{3_1} 展现了本研究的整体架构，它分为 IR 建立、Pass 集合、调优和后端四个部分。这四个功能部件相互关联：IR 建立层负责构建线程块级别和线程级别的中间表示，使算子逻辑能够被结构化描述，并便于 Pass 在不同层级上发挥作用；Pass 集合基于确定的转换规则，先后对计算调度和访存行为实施逐级细化重构；调优部分通过自动探索搜索空间，并将结果反馈给 Pass 集合以产生较优参数配置；后端则面向不同平台完成高性能代码生成与执行。
具体而言，本文的主要贡献包括以下几个方面：
\begin{itemize}
\item 本文提出一种面向硬件层次结构的多级中间表示设计方法。本文在多级表示机制 MLIR 基础上，形成线程块级与线程级一起工作的算子建模方法，让数据分块、并行映射、计算调度策略能够在 IR 层进行清晰表达，为后面的优化及转换提供结构化语义支持。

\item 本文针对 Transformer 关键算子设计并实现了一组优化 Pass，以提高其性能。本文围绕矩阵乘和注意力机制的计算与访存特点，形成包括软件流水线、规约轴并行、线程块重映射、离散访存优化的优化方式，对计算调度和数据访问路径进行整体改进。

\item 本文构建了与 IR 优化流程紧密结合的自动调优策略。本文通过对结构化调优参数空间进行定义，并把剪枝策略同并行编译—测试机制结合起来，实现对不同硬件平台的高效参数搜索，让生成算子在多种异构环境下都可以获得稳定且可迁移的高性能表现。
\end{itemize}

经上述设计后，本文提出的方法可以在 IR 层统一表示算子的结构并对优化策略进行描述，同时能在不同硬件平台之间完成结构复用及性能迁移，这为 Transformer 模型主要算子的自动生成和跨平台优化给出了一套整体的解决方式。接下来的各节将围绕 IR 建模与转换、优化 Pass 设计、自动调优方法进行详细说明。

% ########################################################################################
%                          2 面向硬件层次结构的多级IR设计与转换
% ########################################################################################
\section{面向硬件层次结构的多级IR设计与转换}
本节首先针对传统 GPU 的硬件层次结构，提出线程块和线程层级的算子建模方案，并以多头注意力机制为例进行详细阐述；其次提出一种索引推断算法，用于解决线程块层级 IR 向线程层级 IR 降级过程中的索引映射问题；最后从双层 IR 代码角度出发，通过示例说明 IR 代码的降级过程。


% =============================================================
%                      2.1 算子的计算建模
% =============================================================
\subsection{算子的计算建模}
现代 GPU 等异构加速硬件在其编程模型中通常呈现“线程块-线程”两级并行层级。通过 BlockId 与 ThreadId 构成的张量访存索引表达式，系统能够将计算任务精确映射至对应的线程块及其内部线程。针对这种层次化的硬件特征，本文提出了一种解耦的计算建模策略：将算子的运算逻辑划分为“线程块级”与“线程级”两个层次，并在中间表示的构建中使用不同粒度的 Operation 进行刻画。

\begin{table}[htbp]
\caption{DeepGen 方言中自定义 Operation 列表}\label{tab:table2}
\vspace{0.5em}\centering\wuhao
\renewcommand{\arraystretch}{1.35}
\begin{tabular}{cc}
\toprule[1.5pt]
Operation & 描述 \\ \midrule[1pt]
KernelOp & 描述整个算子计算流程，所有操作均应定义在其内部 \\
ParallelOp & 描述网格的线程块分配及参与计算的线程块内线程数 \\
AllocBufferOp & 创建涵盖共享内存、寄存器及全局内存的缓冲区 \\
CopyOp & 执行数据在共享内存、寄存器与全局内存间的拷贝操作 \\
FillOp & 对寄存器缓冲区进行数据初始化与填充 \\
MatmulOp & Block 级矩阵乘操作，输入共享内存或寄存器，输出至寄存器 \\
BlockOp & Block 级计算，涵盖非矩阵乘计算与特定的数据拷贝逻辑 \\
ReduceOp & 执行包括 max、min、add 等在内的规约（Reduction）操作 \\
EndOp & 自定义操作的块内部终止符，无具体计算含义 \\ \bottomrule[1.5pt]
\end{tabular}
\vspace{\baselineskip}
\end{table}

在线程块层级的逻辑建模中，张量访存的索引计算仅依赖于 BlockId。该层级主要聚焦线程块层级的访存模式与计算调度，屏蔽块内线程的具体任务分配细节。该阶段的索引计算与任务分配策略层次较高，因此其 IR 表达也更为简洁。为了描述算子在线程块级别的计算行为，本文自定义了 MLIR 方言—— DeepGen Dialect。表 \ref{tab:table2} 详细列出了该方言中的核心 Operation 及其语义，线程块层级的 IR 将由这些自定义 Operation 构建。
线程层级的建模则重点转向块内线程的细粒度任务分配与计算实现。该层级通过引入 ThreadId，将线程块级的 Operation 进一步降级为若干个更为低级语义的线程级 Operation，从而确保整个 IR 的逻辑完整性与可执行性。

\begin{table}[htbp]
\caption{自定义缓冲区标识符}\label{tab:table3}
\vspace{0.5em}\centering\wuhao
\renewcommand{\arraystretch}{1.35}
\begin{tabular}{cc}
\toprule[1.5pt]
缓冲区标识符 & 含义说明 \\ \midrule[1pt]
$\mathrm{GM}(X)$ & 缓冲区 X 的地址空间为全局内存 (GM: Global Memory) \\ 
$\mathrm{SM}(X)$/$SmemX$ & 缓冲区 X 的地址空间为共享内存 (SM: Shared Memory) \\ 
$\mathrm{Reg}(X)$/$RegX$ & 缓冲区 X 的地址空间为寄存器 (Reg: Registers) \\ \bottomrule[1.5pt]
\end{tabular}
\vspace{\baselineskip}
\end{table}

除此之外，对于算子中的数据及其对应的内存抽象，本文采用 MLIR 提供的 MemRef 类型（Type）来表示。MemRef 作为 MLIR 的内置方言，提供了数据索引变换的基础设施，便于给出不同内存空间下的数据表达，且转换和优化 Pass 较为丰富。对于线程级细粒度的计算，本文也引入了 MLIR 内置的 Arith、Affine 等方言进行表示。如表 \ref{tab:table3} 所示，为本文中所使用到的内存空间标识符，其中共享内存和寄存器还可以分别使用 $SmemX$ 和 $RegX$ 对缓冲区内存空间进行描述。

\begin{figure}[htbp]
\centering
\includegraphics[width=1\textwidth]{3_2}
\caption{多头注意力机制在 DeepGen 下建模示例的数据流图}\label{3_2}
\vspace{\baselineskip}
\end{figure}

以 Transformer 模型中的多头注意力机制为例，图 \ref{3_2} 展示了其计算过程在双层 IR 模型下的数据流向。图中蓝色区域代表算子在线程块层级的 Operation 表达，橙色区域则为对应 Operation 向线程层级细化的内部结构。各 Operation 采用圆角矩形标识，张量数据缓冲区采用直角矩形标识；数据名称中的 GM、SM、Reg 分别对应全局内存、共享内存及寄存器。涉及到数据访存的 Operation，其下方均标注了相应的索引表达式。图中利用紫色矩形框表示线性循环结构，粉色方框表示非线性循环结构，循环体的迭代变量、上下界及步长均已显式标明。图中所示，其输入为 $\mathrm{GM}(Q)$、$\mathrm{GM}(K)$ 与 $\mathrm{GM}(V)$，输出为 $\mathrm{GM}(O)$，各张量的形状由下方中括号标注（其中 $b$、$h$、$s$、$d$ 分别代表批大小 Batch Size、注意力头数 Head Number、序列长度 Sequence Length 以及注意力头维度 Head Dimension）。红色虚线菱形箭头表示 Operation 之间的数据依赖关系，例如在图 \ref{3_2}(a) 中所示，ReduceOp 必须在 CopyOp 执行完毕后方可启动。

在线程块层级的运算建模中，Operation 访问数据时仅使用 BlockId 参与索引计算。如图 \ref{3_2}(a) 所示，第一个 MatmulOp 所需的输入参数为共享内存中的张量，它们由 CopyOp 从全局内存中拷贝得到。其中 $\{bz, by, bx, 0\}$ 与 $\{bz, by, i, j\}$ 分别表示向全局缓存 $\mathrm{GM}(Q)$ 与 $\mathrm{GM}(K)$ 发起线程块级拷贝的索引表达（$i$、$j$ 为循环变量）。在计算 $Q$ 与 $K$ 的矩阵乘时，为了减少共享内存压力，本文对 $\mathrm{SM}(K)$ 采用切块加载并累加的策略（分块尺寸为 $[16, 32]$），使得 MatmulOp 在单次迭代中仅需针对 $\mathrm{SM}(Q)[j:j+16, 32]$ 与 $\mathrm{SM}(K)[16, 32]$ 的子分块进行计算。

在 DeepGen 方言的设计中，KernelOp 与 ParallelOp 用于封装算子的全局计算域，定义线程块网格分布及块内线程规模，自身不涉及具体的内存访问与算术运算。而 BlockOp 则承担了线程块级别的逐元素（Element-wise）计算与部分数据访存任务。如图 \ref{3_2}(a) 所示，表达式 $acc\_s = exp(acc\_s - m)$ 明确界定了 BlockOp 内部执行的标量/向量计算逻辑，从而在线程块级 IR 中将此类泛型计算进行了统一且高效的封装。

图 \ref{3_2}(b) 进一步展示了部分线程块级 Operation 向线程层级细化的建模过程。在此阶段，线程块级的 Operation 被展开为若干个细粒度的线程级 Operation。虚线指明了 MatmulOp 与 ReduceOp 在线程层级的内部实现细节。

线程层级下，寄存器需要在线程上进一步划分，该层级下的访存、计算等操作均需要引入 ThreadId 进一步细化索引计算。以图 \ref{3_2}(a) 中的 MatmulOp 为例，在降级到线程层级的 IR 表达后，其整体输入输出参数仍与线程块层级的表达一致，即以 $\mathrm{SM}(Q)$ 和 $\mathrm{SM}(K)$ 为输入，$\mathrm{Reg}(acc\_s)$ 为输出。但其内部不再简单地使用 MatmulOp 笼统概括，而是给出由 arith、affine 等细粒度 Operation 组成的具体计算逻辑：共享内存块被加载到寄存器 $\mathrm{Reg}(Q)$ 和 $\mathrm{Reg}(K)$ 中，由线程块中每个线程执行外积计算并累加得到最终结果，再存回到 $\mathrm{Reg}(acc\_s)$ 中。其中，数据从共享内存拷贝进入寄存器时，索引 $\{bk, map1(tid, i)\}$ 与 $\{bk, map2(tid, i)\}$ 分别表示访问 $\mathrm{SM}(Q)$ 和 $\mathrm{SM}(K)$ 的索引表达，$tid$ 表示线程块中的线程索引序号，$i$ 表示每个线程取数的循环变量，map 则为与 $tid$ 和 $i$ 相关的索引重映射函数。

为了更细粒度地表达线程层级的访存、计算等语义，本文采用了 MLIR 中内置的 Affine、Arith 以及 Math 等方言。其中，Affine 方言包含循环变换和仿射映射相关操作，例如循环结构（affine.for）、数据加载/存储操作（affine.load/store）等；Arith 方言包含基本算术操作，例如加、减、乘、除操作（arith.addf/subf/mulf/divf）等；Math 方言则包含数学运算操作，例如指数运算操作（math.expf）、对数运算操作（math.logf）等。以上操作支持本文表达线程层级的复杂计算语义，同时暴露更多优化细节，更好地支持后续优化。

需要说明的是，本文选择矩阵乘法和多头注意力机制作为主要研究对象，是因为二者在 Transformer 推理过程中占据主要计算开销，并且能够较充分地体现分块计算、线程协作、寄存器复用和多级存储访存优化等核心问题。但双层 IR 的设计并不局限于上述两类算子。对于 Transformer 中的其他常见算子，可以依据其计算特征映射到现有 Operation 组合中，从而在相同的线程块级与线程级框架下完成表达和降级。

具体而言，LayerNorm 主要由均值计算、方差计算、归一化和仿射变换组成，其中均值和方差可由 ReduceOp 表达为沿特定维度的求和规约，归一化和缩放平移可由 BlockOp 结合 Arith、Math 等内置方言表达为逐元素计算。GELU 等激活函数属于典型的逐元素非线性算子，其计算不涉及复杂的跨线程数据依赖，可由 BlockOp 表示块级计算语义，并在降级后展开为线程级的 arith 和 math 操作。Softmax 则可分解为行内最大值规约、指数变换、求和规约以及归一化四个阶段，其中最大值与求和阶段由 ReduceOp 建模，指数和除法归一化阶段由 BlockOp 建模。事实上，图 \ref{3_2} 所示的多头注意力机制中已经包含了 Softmax 的主要计算流程，说明该类规约与逐元素混合算子能够被双层 IR 统一描述。

因此，从表示能力上看，DeepGen 方言可以覆盖计算密集型算子、规约型算子以及逐元素算子三类 Transformer 常见计算模式。其扩展过程通常只需要在前端将目标算子的计算过程拆分为 MatmulOp、ReduceOp、BlockOp 和 CopyOp 等基本 Operation 的组合，后续步骤相同。

% =============================================================
%                      2.2 IR降级信息推断算法
% =============================================================
\subsection{IR 降级信息推断算法}

在上一节中，本文将算子建模过程划分为线程块层级与线程层级两个阶段。对应到 MLIR 实现中，首先使用 DeepGen 方言对算子逻辑进行表达，其次通过 IR 降级过程将线程块层级的操作进一步降级为线程层级的操作。在该降级过程中，需要通过引入 ThreadId 参与索引计算，从而将线程块层级的访存操作具体映射到每个线程的执行行为。因此，在执行 IR 降级之前，必须明确 BlockId 与 ThreadId 之间的索引映射关系。

然而，在线程块层级的 IR 表示中，并未显式包含线程层级的索引映射信息，因此无法直接完成线程级访存索引的生成。为解决该问题，需要结合算子的计算语义以及线程块层级 Operation 的结构特性，对线程级索引映射关系进行推断，从而为后续 IR 降级提供必要的指导信息。

\begin{figure}[htbp]
\centering
\includegraphics[width=0.75\textwidth]{3_3}
\caption{LowerInfo 结构体定义}\label{3_3}
\vspace{\baselineskip}
\end{figure}

如图 \ref{3_3} 所示，本文设计了一种名为 LowerInfo（Lowering Information）的结构体，用于记录缓冲区在不同执行层级上的访存索引映射关系。对于每一个涉及缓冲区访问的 Operation，系统都会为其操作的缓冲区建立对应的 LowerInfo 结构。这里的缓冲区在 MLIR 中以 MemRef 类型的 Value 表示。LowerInfo 结构内部采用 MLIR 的 AffineMap 来描述不同层级的索引映射表达式，并通过设置索引变量的上界约束保证访存地址的合法性。

在 LowerInfo 结构中，除了线程块与线程层级的索引映射之外，还引入了线程束层级的索引表达，从而能够实现更细粒度的线程访问控制。此外，结构中还包含用于访存离散化的索引映射机制，通过对线程访问的数据块进行分散化访问，可以有效降低共享内存中的 Bank 冲突问题。针对双层离散化策略的具体实现将在第 \ref{ch3_5} 节中进行详细介绍。最终，在 IR 降级阶段，对于共享内存而言，只需将线程块层级 Operation 中的块级索引表达式与 LowerInfo 中记录的线程、线程束等层级的索引表达式进行组合，而对于寄存器而言只需要组合 LowerInfo 中线程层级的索引表达式，即可得到最终访存索引。

通过对算子执行过程的分析可以发现，不同 Operation 之间对缓冲区索引的依赖关系存在明显差异。例如，在图 \ref{3_2}(a) 中，CopyOp 与 MatmulOp 之间的索引访问彼此独立。具体而言，CopyOp 对共享内存 $\mathrm{SM}(K)$ 进行写入时的索引分布，并不要求与 MatmulOp 从 $\mathrm{SM}(K)$ 读取数据时的索引保持一致。而另一方面，MatmulOp 与 ReduceOp 之间通过寄存器 $\mathrm{Reg}(acc\_s)$ 传递中间结果，因此两者在索引访问上必须保持一致。由于 MatmulOp 计算得到的部分结果分别存储在不同线程的寄存器中，因此在 ReduceOp 执行时，每个线程必须按照 MatmulOp 中数据写入 $\mathrm{Reg}(acc\_s)$ 时的索引方式进行读取，才能保证计算结果的正确性。

基于上述分析可以发现，仅寄存器的 LowerInfo 需要在多个 Operation 之间保持一致，并参与跨 Operation 的依赖关系推断；而对于共享内存或全局内存等其他类型缓冲区，其 LowerInfo 仅在当前 Operation 内部有效。当这些缓冲区在后续 Operation 中再次被访问时，需要根据新的 Operation 语义重新推断其 LowerInfo。因此，本文提出了一种 LowerInfo 推断算法 \ref{alg1}，用于自动生成算子中各访存 Operation 对应的缓冲区索引映射信息，从而指导 IR 的降级过程。根据推断过程中对前序信息依赖程度的不同，线程块层级的 Operation 可分为两类。

\begin{figure}[htbp]
\centering
\includegraphics[width=0.75\textwidth]{3_4}
\caption{MatmulOp 中缓冲区的 LowerInfo 的独立推断}\label{3_4}
\vspace{\baselineskip}
\end{figure}

第一类为可独立推断型 Operation。这类 Operation 可以根据其自身的计算语义直接推导出相关缓冲区的 LowerInfo。例如图 \ref{3_2}(a) 中的 MatmulOp，其输入缓冲区的索引分布可以由输出缓冲区的索引关系推导得到。如图 \ref{3_4} 所示为 MatmulOp 中缓冲区 LowerInfo 的推断过程。其中 $\mathrm{Reg}(X)$ 表示寄存器层级的缓冲区，$Tx$ 表示线程束内部线程的索引，不同颜色表示不同线程束负责计算缓冲区中的不同区域。由于该矩阵乘法采用外积（Outer Product）的计算方式，因此输出寄存器 $\mathrm{Reg}(acc\_s)$ 的 LowerInfo 可以直接根据矩阵乘法的计算规则确定。在计算过程中还引入两个临时寄存器 $\mathrm{Reg}(Q)$ 与 $\mathrm{Reg}(K)$，分别用于存储来自 $\mathrm{SM}(Q^{\mathrm{T}})$ 和 $\mathrm{SM}(K)$ 的数据。尽管这两个寄存器的 LowerInfo 未在图中显式展示，但其索引关系同样可以由 $\mathrm{Reg}(acc\_s)$ 的 LowerInfo 推导得到。因此，通过外积计算的结构特性即可进一步确定 $\mathrm{SM}(Q^{\mathrm{T}})$ 与 $\mathrm{SM}(K)$ 的 LowerInfo 信息。该推断过程对应算法 \ref{alg1} 中的第 8 至第 15 行。

\SetAlgorithmName{算法}{算法}{算法列表}
\SetKwInput{KwIn}{输入}
\SetKwInput{KwOut}{输出}
\SetKwInput{KwData}{输入}
\SetKwInput{KwResult}{输出}
\SetKwFor{For}{for}{do}{end for}
\SetKwFor{ForEach}{for each}{do}{end for}
\SetKwIF{If}{ElseIf}{Else}{if}{then}{else if}{else}{end if}
\SetKwFor{While}{while}{do}{end while}
\DontPrintSemicolon

\SetAlFnt{\xiaosi} % 使用小号字体
\SetAlCapFnt{\xiaosi} % 标题也用小号字体
% \SetKwSty{\xiaosi}

\begin{algorithm}
\caption{LowerInfo 推断算法}\label{alg1}
\setstretch{1.5}
    \KwIn{kernel\_op - 块级别的 KernelOp 操作实例}
    \KwOut{成功返回 (\textbf{op\_bufs\_info\_map}, \textbf{frag\_buf\_info\_map})，失败返回 false}

    // 步骤 0：初始化数据结构 \\
    \textbf{need\_infer\_ops} $\gets$ $\emptyset$ \\
    \textbf{op\_bufs\_info\_map} $\gets$ $\emptyset$ \\
    \textbf{frag\_buf\_info\_map} $\gets$ $\emptyset$ \\

    // 步骤 1：收集需要推断的操作 \\
    \textbf{need\_infer\_ops} $\gets$ \{op $\in$ \textbf{kernel\_op} $|$ op 类型为 CopyOp、BlockOp、MatmulOp、FillOp 或 ReduceOp\} \\

    // 步骤 2：直接推断 \\
    \ForEach{op $\in$ \textbf{need\_infer\_ops}（按前序遍历顺序）}{
        \If{op 满足直接推断条件}{
            \textbf{infos} $\gets$ op.inferLowerInfo() \\
            \textbf{op\_bufs\_info\_map}[op] $\gets$ \textbf{infos} \\
            从 \textbf{infos} 中更新寄存器级 LowerInfo 至 \textbf{frag\_buf\_info\_map} \\
            \textbf{need\_infer\_ops} $\gets$ \textbf{need\_infer\_ops} $-$ \{op\}
        }
    }
    // 步骤 3：迭代式依赖推断 \\
    \For{$k \gets 1$ \KwTo $1000$}{
        \If{\textbf{need\_infer\_ops} = $\emptyset$}{
            \Return (\textbf{op\_bufs\_info\_map}, \textbf{frag\_buf\_info\_map})
        }
        
        \textbf{del\_ops} $\gets$ $\emptyset$ \\
        \ForEach{op $\in$ \textbf{need\_infer\_ops}}{
            \If{基于 \textbf{frag\_buf\_info\_map} 中的已有信息可推断 op}{
                \textbf{infos} $\gets$ 基于依赖关系推断(op, \textbf{frag\_buf\_info\_map}) \\
                \textbf{op\_bufs\_info\_map}[op] $\gets$ \textbf{infos} \\
                从 \textbf{infos} 中更新寄存器级 LowerInfo 至 \textbf{frag\_buf\_info\_map} \\
                \textbf{del\_ops} $\gets$ \textbf{del\_ops} $\cup$ \{op\}
            }
        }
        
        \textbf{need\_infer\_ops} $\gets$ \textbf{need\_infer\_ops} $-$ \textbf{del\_ops}
    }
    
    \Return false
\end{algorithm}

\begin{figure}[htbp]
\centering
\includegraphics[width=1.0\textwidth]{3_5}
\caption{ReduceOp（左）和BlockOp（右）中缓冲区的 LowerInfo 推断}\label{3_5}
\vspace{\baselineskip}
\end{figure}

第二类为依赖型 Operation。这类 Operation 的 LowerInfo 推断需要依赖于前序 Operation 中已经得到的寄存器缓冲区索引信息。例如图 \ref{3_2}(a) 中的 ReduceOp 以及其后的 BlockOp，其缓冲区 LowerInfo 推断过程如图 \ref{3_5} 所示。在 ReduceOp 的推断过程中，输入寄存器 $\mathrm{Reg}(acc\_s)$ 的 LowerInfo 已经在前面的 MatmulOp 阶段推导完成，并存储于 frag\_buf\_info\_map 中。当遍历到 ReduceOp 时，系统结合 $\mathrm{Reg}(acc\_s)$ 的 LowerInfo 以及 ReduceOp 内部的计算逻辑，即可推导得到输出寄存器 $\mathrm{Reg}(m)$ 的 LowerInfo。同理，在后续 BlockOp 中，其计算表达式为 $scale=exp(prev\_m-m)$，在已知 $\mathrm{Reg}(m)$ 索引分布的情况下，结合 Operation 内部的计算规则即可进一步推导出 $\mathrm{Reg}(scale)$ 与 $\mathrm{Reg}(prev\_m)$ 的 LowerInfo 信息。推断完成后，新生成的 LowerInfo 将被存入 frag\_buf\_info\_map 中，以供后续 Operation 继续使用。该过程对应算法 \ref{alg1} 中的第 17 至第 31 行。

从适用范围上看，LowerInfo 推断算法主要面向静态控制流和规则张量访问模式下的块级算子。此类算子的 Operation 拓扑关系在编译期可确定，访存索引可表示为 BlockId、ThreadId、循环变量和仿射映射的组合，且跨 Operation 的数据传递主要通过寄存器缓冲区完成。矩阵乘法、多头注意力机制以及前文讨论的规约型、逐元素型算子均属于该类场景。在这些场景下，算法可以通过前序遍历和迭代式依赖传播逐步补全各 Operation 的 LowerInfo，并在无法继续推断时返回失败，以避免生成不确定的线程级访存映射。

对于包含复杂数据依赖的情形，算法需要满足更严格的前提。若多个 Operation 之间存在分支控制流，则只有当分支条件和两侧缓冲区布局在编译期可判定且汇合点处的寄存器 LowerInfo 一致时，才能继续沿用该推断过程；而对于存在循环携带依赖的场景，若循环每次迭代中的寄存器布局保持不变，且依赖关系可由静态数据流分析确认，则算法仍可将其视为稳定的迭代依赖进行处理；若迭代过程中缓冲区的线程级布局发生变化，或者访存地址依赖运行时数据，则当前算法无法直接完成自动推断。对于无法自动推断的情况，DeepGen 将通过断言终止编译，避免生成错误代码。因此，本文提出的 LowerInfo 推断算法更适合用于结构规则、访存模式明确的 Transformer 关键算子，对动态控制流和数据相关访存的支持仍是后续工作需要进一步扩展的方向。

% =============================================================
%                      2.3 IR降级转换
% =============================================================
\subsection{IR 降级转换}

在上一节中，本文介绍了算子从线程块层级 IR 降级至线程层级 IR 过程中所依赖的 LowerInfo 结构体及其推断方法。
\begin{figure}[htbp]
\centering
\includegraphics[width=0.92\textwidth]{3_6}
\caption{注意力机制在线程块层级 IR}\label{3_6}
\vspace{\baselineskip}
\end{figure}
本节进一步说明 IR 的具体降级转换过程，并以图 \ref{3_2} 所示注意力机制中的 MatmulOp 与 ReduceOp 为示例进行说明。
\begin{figure}[htbp]
\centering
\includegraphics[width=0.92\textwidth]{3_7}
\caption{降级后 MatmulOp 在线程层级的 IR}\label{3_7}
\vspace{\baselineskip}
\end{figure}
图 \ref{3_6} 展示了图 \ref{3_2}(a) 数据流图在 MLIR 中的线程块层级 IR 表示，而图 \ref{3_7} 与图 \ref{3_8} 分别给出了图 \ref{3_6} 中第 15 行和第 20 行代码表示的 MatmulOp 与 ReduceOp 在完成降级后的线程层级 IR 表示。

\begin{figure}[htbp]
\centering
\includegraphics[width=0.92\textwidth]{3_8}
\caption{降级后 ReduceOp 在线程层级的 IR}\label{3_8}
\vspace{\baselineskip}
\end{figure}
这两个 IR 片段同时对应图 \ref{3_2}(b) 左侧与右侧的变换结果。其中 IR 中的 $SmemK$ 与 $SmemQ$ 分别表示 $\mathrm{SM}(K)$ 与 $\mathrm{SM}(Q)$，其余缓冲区名称均与图 \ref{3_2}(a) 中的标注保持一致。

对于示例中的 MatmulOp，其输入缓冲区均位于共享内存，因此 $\mathrm{SM}(Q)$ 与 $\mathrm{SM}(K)$ 的加载索引需要依赖输出寄存器 $\mathrm{Reg}(acc\_s)$ 的计算索引。在上一节中已经说明，$\mathrm{Reg}(acc\_s)$ 的索引分布可以根据 MatmulOp 自身的计算语义直接推导得到，因此该 Operation 涉及的所有缓冲区的 LowerInfo 均可以独立推断得到。如图 \ref{3_9} 所示为 MatmulOp 中各缓冲区对应的 LowerInfo 打印结果。通过结合调优参数信息，例如离散化前后线程块、线程束与线程所处理的数据分块尺寸，可以计算得到各缓冲区对应的布局、索引上界以及索引表达式。

在完成 LowerInfo 推断后，即可执行 MatmulOp 的 IR 降级。如图 \ref{3_7} 所示，降级后的 MatmulOp 在 IR 中通过 affine.load 与 affine.store 操作对 $\mathrm{SM}(Q)$、$\mathrm{SM}(K)$ 以及 $\mathrm{Reg}(acc\_s)$ 进行访存。其中线程访问共享内存时的索引表达式由 LowerInfo 中各层级索引组合得到；而线程访问寄存器内部数据时，则直接使用 thread\_inner\_index 中记录的线程内部索引表达式。因此，在确定所有缓冲区的索引映射关系之后，即可完成 MatmulOp 从线程块层级 IR 到线程层级 IR 的降级转换。

与 MatmulOp 不同，ReduceOp 的 LowerInfo 推断依赖于前序 Operation 的寄存器索引信息。如果在未知 $\mathrm{Reg}(acc\_s)$ 的 LowerInfo 情况下执行推断，将无法直接确定 $\mathrm{Reg}(m)$ 的索引分布。
\begin{figure}[htbp]
\centering
\includegraphics[width=0.85\textwidth]{3_9}
\caption{MatmulOp 中 $\mathrm{SM}(Q)$、$\mathrm{SM}(K)$ 以及 $\mathrm{Reg}(acc\_s)$ 对应的 LowerInfo}\label{3_9}
\vspace{\baselineskip}
\end{figure}
这是因为寄存器在 GPU 执行模型中为线程私有资源，当同一寄存器缓冲区被多个 Operation 访问时，其索引映射必须保持一致，否则将可能导致计算错误甚至访存越界等问题。

如图 \ref{3_10} 所示为 ReduceOp 中缓冲区的 LowerInfo 打印结果，其中注释表示组合后的索引表达式。在推断 $\mathrm{Reg}(m)$ 的 LowerInfo 时，首先利用 ReduceOp 的 $dim$ 属性确定该 Operation 沿行方向执行区间最大值规约操作。同时，根据 $\mathrm{Reg}(acc\_s)$ 的 LowerInfo 可以得知每个线程所持有的数据分块形状为 [2, 4]。在沿行方向执行规约后，每个线程最终仅保留长度为 2 的结果向量，因此 $\mathrm{Reg}(m)$ 的数据分块形状为 [2]。由此可以得到 $\mathrm{Reg}(m)$ 的 LowerInfo，其索引分布等价于 $\mathrm{Reg}(acc\_s)$ 在行方向上的映射拷贝。

\begin{figure}[htbp]
\centering
\includegraphics[width=0.85\textwidth]{3_10}
\caption{ReduceOp 中 $\mathrm{Reg}(acc\_s)$ 和 $\mathrm{Reg}(m)$ 对应的 LowerInfo}\label{3_10}
\vspace{\baselineskip}
\end{figure}

如图 \ref{3_8} 所示为 ReduceOp 降级后的线程层级 IR。其中第一步对 $\mathrm{Reg}(m)$ 缓存执行清空操作，以保证数据正确性；随后使用 $\mathrm{Reg}(acc\_s)$ 进行最大值规约计算，并将结果写入 $\mathrm{Reg}(m)$，其索引直接使用 thread\_inner\_index 中存储的表达式。之后根据 LowerInfo 中的布局信息，可以得知整个线程块在 $dim=1$ 方向上的线程束数量和线程数量，并利用 GPU 方言提供的线程束级 Shuffle XOR Operation，实现线程块布局中同一行线程之间的规约操作，最后将结果存储在 $\mathrm{Reg}(m)$ 中。



% ########################################################################################
%                            3 关键算子的高性能优化pass集合
% ########################################################################################
\section{关键算子的高性能优化 Pass 集合}
上一节已经设计了用于表达算子计算逻辑的双层 IR，本节重点讨论算子的优化方案。在 MLIR 机制中，Pass 是对 IR 实施整体程序变换的核心方式。每个 Pass 基于模式匹配和重写规则，能够完成操作替换、方言降级、控制流与数据流重构、硬件映射标注等操作。DeepGen 面向 Transformer 架构中耗时较高的关键算子，建立了一套专门用于优化的 Pass 机制。这些 Pass 主要聚焦于两类性能瓶颈：一类是计算调度不足导致的线程占用率不高和算术强度不足，另一类是访存模式不合理导致的全局内存带宽浪费、共享内存 Bank 冲突以及 L2 Cache 命中率不高。本文将优化 Pass 分为计算调度优化和访存层次优化两类，两者相互配合，以提升端到端执行性能。

计算调度优化 Pass 负责循环结构的切分、重排、并行化、软件流水、规约维度拆分等变换，用于提高线程级与块级并行度，并增强计算与访存的重叠能力。访存优化 Pass 则通过全局内存合并访存、共享内存多缓冲、动态分配、Bank 冲突消除、L2 Cache 命中率提升、向量化加载等手段，使算子在实际硬件上尽可能接近理论带宽与理论计算屋顶线（Roofline Line）。表 \ref{tab:table4} 列出了 DeepGen 引入的 13 项核心优化 Pass 及其主要功能。


\begin{table}[htbp]
\caption{算子优化 Pass 集合列表}\label{tab:table4}
\vspace{0.5em}\centering\wuhao
\renewcommand{\arraystretch}{1.35}
\begin{tabularx}{\textwidth}{c >{\centering\arraybackslash}X}
\toprule[1.5pt]
Pass 名称 & 描述 \\ \midrule[1pt]
Split & 循环切分，将单个循环按照循环边界进行切分得到多个循环 \\
Parallelize & 循环并行化，将串行循环转换成并行执行 \\
Unroll & 循环展开，按照展开次数将循环进行展开 \\
SoftwarePipeline & 软件流水线优化，将循环体内符合“访存-计算”的循环执行软件流水线操作 \\
Schedule & Operation 调度，通过调整 IR 中 Operation 的位置从而调度程序语句的执行次序 \\
LocalSplitK & 局部规约轴切分，将 MatmulOp 中规约循环切分为循环后并进行并行化 \\
CoalescedAccess & 全局访存合并，检测缓冲区的行列连续性，保证对于全局内存访问连续性 \\ 
MultiBuffer & 翻倍扩大缓冲区大小，并修改访存索引，为软件流水线提供缓冲区基础 \\ 
BlockSwizzle & 线程块索引重映射，通过修改线程块 ID 索引的映射，优化 L2 Cache 命中率 \\ 
Vectorize & 向量化访存，识别访存操作并向量化访存减少 IO 指令个数，提高指令级并行 \\ 
AccessScatter & 访存离散化，通过离散化线程对于共享内存的访问地址，减少 Bank 冲突 \\ 
MemoryReuse & 缓冲区重用，通过检测缓冲区生命周期，进行缓冲区复用 \\ 
SharedDynamic & 动态化共享内存，将算子中的所有共享内存合并，并修改为动态共享内存 \\ \bottomrule[1.5pt]
\end{tabularx}
\vspace{\baselineskip}
\end{table}


% =============================================================
%                3.1 矩阵乘法与MHA优化分析
% =============================================================
\subsection{矩阵乘法与注意力机制优化分析}
Transformer 进行推理时，主要计算单元包括矩阵乘法和注意力机制。二者在结构形式方面存在差异，但从计算模式和性能特性角度分析，其优化问题表现出较高一致性。本文主要从计算结构、访存行为、并行映射和规约特性等方面对两类算子进行分析，为后续优化设计提供理论依据。

图 \ref{3_11}(a) 从计算结构角度展示了矩阵乘法的三重循环规约方式。矩阵乘法拥有一个重要的规约维度，也就是 K 维。输出矩阵所包括的各个子块在空间上各自独立，输入数据则具有较大的数据复用潜力。如图 \ref{3_11}(b) 所示，在理想情况下，每个输出子块的计算可以在片上存储中完成大部分规约操作，并在有需要时才访问全局内存。在程序运行过程中，由于片上存储空间比较有限，数据需要从全局内存进行多次加载。如果分块策略或者线程映射方式设计不够合理，带宽会成为主要的性能瓶颈。因此，矩阵乘法的理论算术强度虽然比较高，但是它的实际性能经常会被访存效率所限制。
\begin{figure}[htbp]
\centering
\includegraphics[width=0.95\textwidth]{3_11}
\caption{矩阵乘法在 GPU 上的两种计算方式}\label{3_11}
\vspace{\baselineskip}
\end{figure}

注意力机制由两个矩阵乘算子和一个逐行归一化过程组合而成。首先通过查询矩阵和键矩阵相乘来计算相关性矩阵，然后对每一行进行缩放和 Softmax 归一化，最后把归一化结果与值矩阵相乘以得到输出。从计算本质来看，多头注意力机制的主要计算负荷仍然集中在矩阵乘阶段。但是，和基础矩阵乘法计算相比，多头注意力机制在中间处理时会生成相关性矩阵，其规模随序列长度平方增长，这使计算负担和访存需求更加突出。如果中间结果频繁写回全局内存后又被读取，将明显降低整体执行效率。多头注意力机制的优化不仅要考虑计算调度问题，还需要关注中间数据的存储和复用方式。
\begin{equation}\label{softmax}
y_i = \frac{\exp\left(x_i-\max_{1\le k\le V} x_k\right)}
{\sum_{j=1}^{V}\exp\left(x_j-\max_{1\le k\le V} x_k\right)}
\end{equation}
\begin{figure}[htbp]
\centering
\includegraphics[width=0.95\textwidth]{3_12}
\caption{无 HBM 写回的融合多头注意力机制计算流程图}\label{3_12}
\vspace{\baselineskip}
\end{figure}
如图 \ref{3_12} 所示，本文使用 Safe Softmax（公式 \ref{softmax}）来实现多头注意力机制的融合；对比传统多头注意力机制，该方法不需要把 $QK^T$ 的结果写回 HBM，从而减少访存压力。

从性能瓶颈角度看，矩阵乘和多头注意力机制在计算能力充足时都会面临明显的带宽压力。屋顶模型（Roofline Model）表明，当数据复用程度较低时，算子执行会处于带宽受限区域。因此，增加数据在片上存储中的停留时间并提高复用程度，是优化的主要方向。在线程块级执行中，多个线程块经常访问重叠的输入数据。如果线程块以简单的行主序遍历输出空间，时间局部性较弱，缓存命中率会下降。通过调整线程块遍历次序，可使访问相邻数据区域的计算块在时间上依次执行，从而提高缓存利用率并减少全局内存访问。

矩阵乘的规约维度通常在单个线程块内部完成，但当规约维度较大时，单个线程块的计算可能成为性能限制。此时可将规约轴划分为多个子段，由多个线程块并行计算部分结果，再通过额外合并得到最终输出。该策略能够提高并行度，但也会引入归并开销，因此需要平衡并行效率和同步成本。多头注意力机制中的矩阵乘阶段也具有类似特性，因此规约轴并行具有通用意义。此外，计算与访存的执行关系同样会影响性能。传统实现中，数据加载和计算通常串行执行，而全局内存访问延迟远高于算术操作延迟，容易造成计算单元空闲。通过流水线技术重叠数据加载与计算过程，可以隐藏访存延迟并提高整体吞吐量。


% =============================================================
%                3.2 Software Pipeline：软件流水线优化
% =============================================================
\subsection{Software Pipeline：软件流水线优化}

在深度学习算子里，循环结构通常表现为加载、计算和存储，也就是所谓的 load-compute-store 执行模式。在该模式下，计算指令大多依靠前序访存指令提供的数据，进而建立起严格的数据依赖链。
\begin{figure}[htbp]
\centering
\includegraphics[width=1\textwidth]{3_13}
\caption{MHA 中 $QK^T$ 在线程块层级执行软件流水线的变换过程示意图}\label{3_13}
\vspace{\baselineskip}
\end{figure}
在访存延迟比较高的情况下，GPU 需要等内存访问完成以后才可以继续进行计算，这会使计算单元在等待时处于空闲，严重限制了算子性能的发挥。现代 GPU 具备硬件指令流水线、部分乱序发射能力，跨指令数据依赖会使访存和计算阶段进行串行执行，并且使计算核心利用率出现下降。

\begin{figure}[htbp]
\centering
\includegraphics[width=1\textwidth]{3_15}
\caption{MHA 中 $QK^T$ 在软件流水线变换前的线程块层级 IR}\label{3_15}
\vspace{\baselineskip}
\end{figure}
\begin{figure}[htbp]
\centering
\includegraphics[width=1\textwidth]{3_16}
\caption{MHA 中 $QK^T$ 在软件流水线变换后的线程块层级 IR}\label{3_16}
\vspace{\baselineskip}
\end{figure}
为隐藏全局内存访问以及共享内存到寄存器之间的数据访问延迟，本文设计了针对线程块级和线程级的双层软件流水线优化 Pass。该 Pass 首先通过数据流分析识别循环体内 load-compute-store 三阶段结构中严格的前后依赖关系，随后通过双或多缓冲索引重写，将原本单次迭代访存与计算解耦为相互独立的流水阶段，自动构造标准的 Prologue、Main Loop 和 Epilogue 三段式软件流水线。变换后，下一迭代的数据加载可与当前迭代的计算完全重叠，实现访存延迟的有效隐藏，显著提升指令级并行度与实际算子吞吐。
\begin{figure}[htbp]
\centering
\includegraphics[width=1\textwidth]{3_14}
\caption{MHA 中 $QK^T$ 在线程层级执行软件流水线的变换过程示意图}\label{3_14}
\vspace{\baselineskip}
\end{figure}

从图 \ref{3_13} 所示示例可以看出，示例流水线深度为 2，因此 $\mathrm{SM}(K)$ 被划分为两个独立缓冲区。随后 SoftwarePipeline Pass 对线程块层级 IR 进行重排，图 \ref{3_15} 和图 \ref{3_16} 分别展示了执行该 Pass 前后线程块层级 IR 的变化：在 Prologue 阶段提前加载 K 矩阵的第一部分数据；在 Main Loop 阶段，一边加载下一迭代的 K 矩阵数据，一边执行当前迭代的 MatmulOp 计算；最后在 Epilogue 阶段完成剩余未计算的数据块，从而形成完整的软件流水线执行结构。

\begin{figure}[htbp]
\centering
\includegraphics[width=1\textwidth]{3_17}
\caption{MHA 中 $QK^T$ 在软件流水线变换前的线程层级 IR}\label{3_17}
\vspace{\baselineskip}
\end{figure}

在 MatmulOp 内部的线程层级 IR 中，同样存在类似的访存—计算依赖关系。具体而言，在循环中需要先从共享内存 $\mathrm{SM}(Q)$ 和 $\mathrm{SM}(K)$ 中加载数据至寄存器 $\mathrm{Reg}(Q)$ 与 $\mathrm{Reg}(K)$，随后在内层循环中执行外积计算。与线程块层级优化方式类似，首先使用 MultiBuffer Pass 对寄存器缓冲区 $\mathrm{Reg}(Q)$ 与 $\mathrm{Reg}(K)$ 进行扩展，从而形成双缓冲结构，为流水线执行提供数据缓冲基础。图 \ref{3_14} 展示了 MatmulOp 在线程层级应用软件流水线优化的执行流程，而图 \ref{3_17} 和图 \ref{3_18} 则展示了在执行该 Pass 前后线程层级的 IR 变化，在此基础上通过 SoftwarePipeline Pass 对线程层级循环结构进行重排，使寄存器数据加载与外积计算能够在不同迭代阶段并行执行。

软件流水线优化的原理主要是消除访存操作和计算操作之间存在着的直接依赖关系，也就是让计算阶段的输入来源于前一阶段已完成加载的数据，同时当前阶段会为正在进行的下一次迭代准备数据。采用这种方式，在进行第 $n$ 次迭代计算时，系统可以同时载入第 $n+1$ 次迭代所需数据，从而实现访存与计算的重叠执行。在具体进行实现的过程中，原始的循环结构被拆分成 Prologue、Main Loop 和 Epilogue 这三个阶段。其中，Prologue 阶段负责完成初始数据预取。
\begin{figure}[htbp]
\centering
\includegraphics[width=1\textwidth]{3_18}
\caption{MHA 中 $QK^T$ 在软件流水线变换后的线程层级 IR}\label{3_18}
\vspace{\baselineskip}
\end{figure}
Main Loop 阶段是流水线在稳态进行执行的过程中，每次循环迭代都会进行访存操作和计算操作。Epilogue 阶段主要用来进行最后一轮尚未完成的计算任务。主循环将循环起始位置调整到原循环的第二次迭代，循环上界保存不变。访存操作和计算操作使用各不相同的索引偏移，单次迭代中的访存和计算便不会有直接依赖关系，可以并行执行，进而提高 GPU 执行单元的利用率。



% =============================================================
%                3.3 Local Split K：外积规约轴并行优化
% =============================================================
\subsection{Local Split K：外积规约轴并行优化}
在 Transformer 模型中，当矩阵乘法运算的规约轴维度显著大于其他两个维度时，严重限制了算子的并行度，如图 \ref{3_19}(a) 所示。在极端情况下，当输出矩阵的尺寸，即 M 和 N 维度也很小时，每个线程块需要处理的计算任务进一步减少，导致线程块级别的并行度急剧下降，难以充分利用现代 GPU 的大规模并行计算能力。为缓解此问题，一种直观的思路是减小每个线程块的计算负载，如图 \ref{3_19}(b) 所示。然而，这种方法会直接导致线程级别的并行度降低，同样无法有效提升硬件利用率。
\begin{figure}[htbp]
\centering
\includegraphics[width=0.85\textwidth]{3_19}
\caption{长规约轴矩阵乘法的三种计算方式}\label{3_19}
\vspace{\baselineskip}
\end{figure}

针对上述问题，本文设计并实现了 LocalSplitK Pass，通过对规约轴方向的计算任务进行重新划分来提升并行度。其核心思想是将原本由单个线程块完成的规约计算在逻辑上拆分为多个子任务，并在同一线程块内并行执行。如图 \ref{3_19}(c) 所示，原本由一个线程块计算的输出区域被划分为多个独立子区域，同时线程块内部的线程数量按比例扩展。
\begin{figure}[htbp]
\centering
\includegraphics[width=0.75\textwidth]{3_20}
\caption{MatmulOp 操作执行 LocalSplitK Pass 优化过程}\label{3_20}
\vspace{\baselineskip}
\end{figure}
每个线程组负责计算规约轴上不同子区间的部分结果，并将中间结果暂存于共享内存中。计算完成后，通过一次同步操作对共享内存中的部分和进行规约累加，从而得到最终输出结果。

从实现角度来看，该 Pass 通过调整 MatmulOp 外积计算的任务分配策略来挖掘更多线程级并行性。具体而言，其将外积计算中的规约轴维度划分为多个子区间，并将每个子区间分配给线程块内部不同的线程组执行，如图 \ref{3_20} 所示。同时，线程块规模按照切分倍数进行扩展，使每个物理线程承担的计算负载与原始设计保持一致，从而在不改变整体计算语义的前提下，实现规约轴方向的子任务级并行。
\begin{figure}[htbp]
\centering
\includegraphics[width=1\textwidth]{3_21}
\caption{未应用 LocalSplitK Pass 时 MatmulOp 的线程层级 IR}\label{3_21}
\vspace{\baselineskip}
\end{figure}

图 \ref{3_21} 和图 \ref{3_22} 展示了图 \ref{3_20} 所示的 MatmulOp 操作在应用 LocalSplitK Pass 前后的 IR 变化情况，其中线程块中的线程数量为 128。图 \ref{3_21} 给出了未应用优化时 MatmulOp 在线程层级的原始 IR 结构，而图 \ref{3_22} 则为应用 LocalSplitK Pass 后的优化 IR。从两者对比可以看出，该 Pass 将原本完整的外积循环逻辑划分为多个子部分。具体而言，在图 \ref{3_22} 第 7 行的 affine.loadOp 操作中，规约轴方向的索引访问被重构为基于线程组的偏移访问形式，从而实现线程组之间的并行计算，而其余索引则通过 map 函数进行简化表示。

此外，为了整合各线程组的计算结果，该 Pass 在 IR 中引入新的共享内存缓冲区 $smemSP$ 与寄存器缓冲区 $regSP$，用于中间结果的存储与规约。如图 \ref{3_22} 第 26 至第 31 行所示，计算过程中各线程组首先将局部结果写入共享内存 $smemSP$，随后在第 34 至第 45 行通过寄存器 $regSP$ 对共享内存中的数据进行逐步累加，从而恢复完整的外积计算结果。通过这种方式，在保持计算语义不变的前提下实现了规约轴方向的并行化处理。
\begin{figure}[htbp]
\centering
\includegraphics[width=1\textwidth]{3_22}
\caption{应用 LocalSplitK Pass 后 MatmulOp 的线程层级 IR}\label{3_22}
\vspace{\baselineskip}
\end{figure}
最后，在执行 MemoryReuse Pass 后，$smemSP$ 将重用共享内存 $smemA$ 和 $smemB$ 的缓冲区，寄存器 $regSP$ 也会重用 $regC$ 的缓冲区，从而减少因过度使用缓冲区造成内存溢出的问题。


% =============================================================
%                3.4 Block Swizzle：线程块访存重排优化
% =============================================================
\subsection{Block Swizzle：线程块访存重排优化}

在异构计算系统中，减轻内存墙问题是提高整体计算性能的重要目标。软件管理的共享内存能够通过显式局部性管理来控制规则数据访问，但在复杂访存模式中，硬件管理的通用缓存，如 L1 或 L2 Cache ，仍然具有不可替代的作用。通用缓存可以自动聚合线程束内的非合并访问请求，提高带宽利用率，并为指令流和随机访存提供透明的局部性保障。在 Transformer 推理等应用场景中，当连续调度的线程块访问的全局内存区域重叠度较高时，L2 Cache 复用率会随之提高，由此可以减轻全局内存访问压力。

根据上述观察，本文实现了 BlockSwizzle Pass，通过对线程块 ID 进行轻量级重映射来优化线程块调度顺序。该方法建立逻辑 BlockId 和物理 BlockId 之间的映射关系，将默认的行优先分配模式改为分组优先模式，使访问相邻数据区域的线程块在时间上更加靠近，从而改善线程块之间的时空局部性，提高缓存复用率。

\begin{figure}[htbp]
\centering
\includegraphics[width=0.7\textwidth]{3_23}
\caption{线程块层面矩阵乘法的两种不同计算次序方式}\label{3_23}
\vspace{\baselineskip}
\end{figure}

图 \ref{3_23} 以矩阵乘法计算为例展示了 Block Swizzle Pass 的优化原理。图中 A 和 B 为输入矩阵，C 为输出矩阵，每个小方格表示一个线程块负责计算的输出区域，因此计算矩阵 C 需要启动 64 个线程块。如图 \ref{3_23}(a) 所示，在默认的行优先执行模式下，线程块按照行主序依次计算矩阵 C，需要加载 A 的一行数据和整个 B 矩阵的数据，共 72 个数据分块。
\begin{figure}[htbp]
\centering
\includegraphics[width=0.95\textwidth]{3_24}
\caption{BlockSwizzle Pass 作用前的行优先映射}\label{3_24}
\vspace{\baselineskip}
\end{figure}
该计算模式下，缓存中的数据会很快被刷新，当下次需要用到相同数据时，只能从全局内存中重新加载。采用 BlockSwizzle Pass 优化后，连续的四个线程块负责计算如图 \ref{3_23}(b) 矩阵中的阴影所示区域。此场景下仅需读取 48 个数据分块，且线程块访问 A 和 B 矩阵时的数据重叠性更高，缓存中数据刷新速率较低，因此可有效提高缓存命中率。
\begin{figure}[htbp]
\centering
\includegraphics[width=0.95\textwidth]{3_25}
\caption{BlockSwizzle Pass 作用后的分组优先映射}\label{3_25}
\vspace{\baselineskip}
\end{figure}

图 \ref{3_24} 和图 \ref{3_25} 展示了本文实现的两种线程块 ID 映射方式在 IR 层面的表示形式。其中图 \ref{3_24} 为默认的行优先映射方式，而图 \ref{3_25} 则为经过 BlockSwizzle Pass 优化后的分组列优先映射方式，两者分别对应图 \ref{3_23}(a) 与图 \ref{3_23}(b) 中的计算模式。

BlockSwizzle Pass 优化时会识别线程块索引操作（gpu.block\_id），然后采用 AffineMap 记录重映射表达式，并使用 AffineApplyOp 操作执行 AffineMap 中记录的映射方式，从而实现线程块 ID 的重映射过程。在图 \ref{3_23}(b) 中，本文通过调优参数 group\_size\_m 确定分组大小，组内线程块采用列优先方式进行计算。在实际映射过程中，需要考虑计算 C 矩阵时列方向线程块数量能否被 group\_size\_m 整除。因此，如图 \ref{3_23}(b) 所示，本文使用 AffineMinOp 获取每个线程块所在组的实际 group\_size\_m 大小，保证线程块 ID 重映射的正确性。通过该方式对线程块索引进行重映射，能够以较小开销实现算子性能提升。

% =============================================================
%                3.5 Access Scatter：离散化访存优化
% =============================================================
\subsection{Access Scatter：离散化访存优化}\label{ch3_5}

GPU 体系结构将共享内存都被组织为 32 个 Bank，每个 Bank 的宽度为 4 字节。若一个线程束中每个线程请求的数据总量超过 4B，硬件会将其拆分为多个内存事务，每个内存事务为 128B。Bank 冲突的计算以每个事务为单位，即当同一事务内的线程访问同一 Bank 的不同地址时，会发生 Bank 冲突，导致串行访问。冲突路数越多，有效带宽越低，性能下降越严重。因此，优化数据布局以规避事务内的 Bank 冲突至关重要。

\begin{figure}[htbp]
\centering
\includegraphics[width=0.95\textwidth]{3_26}
\caption{MatmulOp 操作无离散化优化的计算布局}\label{3_26}
\vspace{\baselineskip}
\end{figure}

以图 \ref{3_26} 与图 \ref{3_27} 为例，展示了图 \ref{3_2} 中第二个 MatmulOp 操作在不同访存策略下的计算示意。该算子的输入参数 $\mathrm{Reg}(acc\_s)$ 与 $\mathrm{SM}(V)$ 的形状分别为 [32, 32] 和 [32, 128]，分别位于寄存器与共享内存中。由于 $\mathrm{Reg}(acc\_s)$ 位于寄存器中，其线程级数据布局已确定，因此 $\mathrm{SM}(V)$ 与结果 $\mathrm{Reg}(acc\_o)$ 的 LowerInfo 需要根据 $\mathrm{Reg}(acc\_s)$ 的布局进行推导。在该配置下，每个线程持有 $\mathrm{Reg}(acc\_s)$ 中的 [2,4] 数据块，因此需要从 $\mathrm{SM}(V)$ 中加载大小为 [2,16] 的数据块以完成外积计算。

\begin{figure}[htbp]
\centering
\includegraphics[width=0.95\textwidth]{3_27}
\caption{MatmulOp 操作有离散化优化的计算布局}\label{3_27}
\vspace{\baselineskip}
\end{figure}

如图 \ref{3_26} 所示，线程 T0 需要从 $\mathrm{SM}(V)$ 中加载一个 [2,16] 的数据块。由于线程最大向量化加载宽度为 16 字节，每次最多读取 4 个元素，因此该数据块需要通过 4 次迭代完成加载。在每次迭代中，线程需要加载一行 16 个元素，而每行加载通过 4 条 LDS.128 指令完成，对应图中不同灰度标识的数据块。根据 $\mathrm{SM}(V)$ 在共享内存 Bank 上的映射关系可以观察到，在同一迭代中，0-7 号线程会访问同一 Bank 的不同地址，因此产生 4 路 Bank 冲突，从而降低访存效率。

针对这一问题，本文在缓冲区 LowerInfo 结构体中引入离散化索引变量（xxx\_scatter\_index），用于对线程的访存数据进行重新划分与分散索引，使得不同线程访问的数据能够分布到不同 Bank 中。图 \ref{3_26} 与图 \ref{3_27} 分别展示了离散化优化前后计算的访存模式。将访存进行离散化后，原本每个线程需要取连续的 16 个元素，变为每个线程分 4 次分散取 4 个元素；从 $\mathrm{SM}(V)$ 数据到 Bank 的映射关系可以看出，0-7 号线程内部的 4 路 Bank 冲突被消除。

\begin{figure}[htbp]
\centering
\includegraphics[width=1\textwidth]{3_28}
\caption{AccessScatter Pass 优化前 $\mathrm{SM}(V)$ 缓冲区的 LowerInfo}\label{3_28}
\vspace{\baselineskip}
\end{figure}

\begin{figure}[htbp]
\centering
\includegraphics[width=1\textwidth]{3_29}
\caption{AccessScatter Pass 优化后 $\mathrm{SM}(V)$ 缓冲区的 LowerInfo}\label{3_29}
\vspace{\baselineskip}
\end{figure}

该优化 Pass 通过修改线程级别的访存模式实现，即调整 LowerInfo 中与离散化相关的参数。若 xxx\_scatter\_bound 参数的各维上界均为 1，则表示未启用访存离散化优化，此时线程访存索引不会被分散。为实现访存离散化，本文缩小原本线程块、线程束和线程访问的数据分块大小，并使用相同的索引方式进行多次访存；每次访存的偏移由 LowerInfo 中 xxx\_scatter\_index 变量存储的索引表达式计算得到，从而使线程每次访问不同地址，最终实现访存离散化优化。

图 \ref{3_28} 和图 \ref{3_29} 分别展示了 MatmulOp 操作中 $\mathrm{SM}(V)$ 缓冲区在执行 AccessScatter Pass 前后的 LowerInfo 配置，分别与图 \ref{3_26} 和图 \ref{3_27} 中的计算示意相对应。其中，block\_scatter\_index 与 warp\_scatter\_index 分别从线程块与线程束两个层级控制访存离散化模式。通过不同索引组合方式，可以生成多种访存布局，从而为自动调优系统提供更加丰富的优化空间。

综上所述，AccessScatter Pass 通过在 IR 层引入离散化访存索引机制，对线程级共享内存访问模式进行重排，从而有效降低甚至消除 Bank 冲突问题。该方法在不改变计算语义的前提下显著提升共享内存带宽利用率，并为不同算子配置提供灵活的访存布局选择，为后续自动调优过程提供了更多潜在的高性能实现方案。

% =============================================================
%                3.6 硬件平台差异与移植策略
% =============================================================
\subsection{硬件平台差异与移植策略}

上述优化 Pass 虽然在 IR 层采用统一的变换形式，但其最终效果会受到目标硬件平台执行模型和存储层次结构的影响。NVIDIA GPU 与海光 DCU 均采用大规模并行执行模型，但在线程束组织、片上存储结构、缓存层次以及指令调度策略等方面仍存在差异。若忽略这些差异，直接复用单一平台上的参数配置，可能导致共享内存利用率下降、线程占用率不足或缓存复用效果减弱。因此，DeepGen 在跨平台生成代码时，需要将统一 IR 表达与平台相关参数约束结合起来，由自动调优框架为不同硬件选择合适的优化组合。

首先，片上存储和缓存层次差异会影响访存类 Pass 的有效性。对于共享内存或 LDS 这类软件可管理存储，不同平台在容量、Bank 数量、Bank 宽度以及访存合并规则上并不完全一致，因此 AccessScatter Pass 中离散化粒度和索引重排方式需要根据目标平台重新选择。当 Bank 冲突较严重时，离散化访存能够显著提高共享内存带宽；而当目标平台的访存合并机制已经能够较好处理连续访问时，过度离散化可能增加指令数量和索引计算开销。BlockSwizzle Pass 则主要依赖 Cache 的容量、替换策略和线程块调度顺序。当连续线程块之间存在较高数据复用时，线程块重映射能够提高缓存命中率；但若输入形状导致线程块之间数据复用较弱，则该 Pass 的收益会相应降低。

其次，线程束调度和执行粒度差异会影响计算调度类 Pass 的收益。NVIDIA GPU 通常以线程束为基本调度单位，而海光 DCU 的底层执行粒度和线程组织方式与其并不完全相同。SoftwarePipeline Pass 的收益取决于访存延迟、寄存器资源和线程块占用率之间的平衡。较深的流水线能够提高访存与计算的重叠程度，但也会增加寄存器和共享内存占用，从而降低并发线程块数量。LocalSplitK Pass 适合规约轴较长且块级并行度不足的场景，它通过在同一线程块内增加规约子任务提高并行度，但也会引入额外的同步和片上规约开销。因此，不同硬件平台上同一 Pass 的最优参数并不固定，需要结合线程束大小、寄存器数量、共享内存容量和最大线程块规模共同确定。

基于上述分析，DeepGen 移植到新硬件平台时通常需要完成以下工作。第一，补充目标平台的硬件资源描述，包括线程束或执行波前大小、单线程块最大线程数、共享内存或 LDS 容量、寄存器限制、缓存层次和向量化访存宽度等。第二，在后端代码生成层适配目标平台的内核启动接口、地址空间标识和线程索引获取方式，使线程块级 IR 能够正确降级到目标平台的执行模型。第三，根据硬件资源约束调整优化 Pass 的合法性检查和参数范围，例如流水线阶数、LocalSplitK 切分数、BlockSwizzle 分组大小以及 AccessScatter 离散化粒度。第四，利用自动调优框架在目标平台上重新搜索候选配置，并通过数值精度验证和性能测试筛选最优实现。通过这种方式，DeepGen 可以在保持 IR 和优化 Pass 框架基本一致的前提下，将平台相关差异集中到硬件描述、后端适配和调优参数选择中，从而降低迁移到新平台时的工程成本。

% ########################################################################################
%                            4 基于空间剪枝的自动调优框架设计
% ########################################################################################
\section{基于空间剪枝的自动调优框架设计}
本文在前面所述的 IR 优化方法基础上，形成了一种利用空间剪枝进行自动调优的策略。该策略在 IR 层定义结构化调优参数方法，对算子执行结构进行参数化描述，并在此基础上形成可搜索的调优空间。本文使用若干与硬件资源约束和算子结构特性相关的剪枝策略，在调优阶段事先剔除不合理或低效的参数组合，从而缩小搜索空间并提高调优效率。通过自动编译、运行和性能评估，该策略可以在合理时间内为目标算子生成高性能实现。本文先介绍自动调优框架的整体设计和工作流程，接着给出调优空间参数的定义及其生成方式，最后对调优空间的剪枝策略做详细说明。

% =============================================================
%                       4.1 调优框架的设计
% =============================================================
\subsection{调优框架的设计}
图 \ref{3_30} 展示了自动调优器 DeepGen 的整体架构及其模块组成。整个调优器框架分为三大组成部分：Compiler、Tester 和 Monitor。Compiler 主要负责算子生成，包括调优空间生成和剪枝模块、代码生成模块以及 Python 加载与封装器；Tester 分为数据收集器和性能测量模块，主要负责收集算子文件与运行时数据，并对算子进行性能测量；Monitor 端提供运行管理，主要包括远程任务管理器和任务控制器，负责控制 Compiler 和 Tester 端的启动、通信与任务同步，保证调优过程顺利执行。自动调优器框架接受用户输入的任务配置文件和调优空间配置文件，输出算子的最优代码生成方案、运行时元数据及其二进制文件。
\begin{figure}[htbp]
\centering
\includegraphics[width=0.85\textwidth]{3_30}
\caption{DeepGen 调优模块总体架构}\label{3_30}
\vspace{\baselineskip}
\end{figure}

在框架运行时，Compiler、Tester 和 Monitor 彼此为独立的进程，可分布在不同的主机上，通过套接字（Socket）进行同步。Monitor 将任务下发给 Compiler 和 Tester 后，Compiler 和 Tester 协同工作，分别进行算子的编译和性能测试。Compiler 和 Tester 之间通过 Monitor 给出的通信端口号进行任务同步。当 Compiler 完成所有代码生成任务后会向 Tester 端发起结束请求。Tester 接收到此请求，完成所有测试任务后发出响应，并将结果回传 Monitor，随后 Compiler 和 Tester 进程结束。Monitor 接收到结果文件后也结束进程。至此，调优过程完成。

基于 Compiler、Monitor 和 Tester 的调优框架设计可以将编译和测试分离，使两者同步进行，下一轮的算子编译不必再等待上一批算子测试完成，测试过程也能充分与编译时间重叠，因此系统总体的速度得到了提高。同时该设计也避免了编译进程过多带来的 CPU 高负载对性能测量结果准确度的影响。

图 \ref{3_31} 展示了 DeepGen 自动调优框架的详细工作流程：首先，DeepGen 对剪枝之后的调优空间进行读取。编译进程池会根据设定的容量大小，开启若干个编译进程，同时编译多个方案。同时，调优框架会对函数签名进行分析，生成算子启动程序。之后，进程池将多个方案编译完毕，生成各自的算子二进制文件，并对其运行时配置进行打包处理，产出 pkl 文件。

接下来，自动调优框架进入调优阶段：对于生成的算子文件，系统根据其 pkl 文件包含的运行时信息，将其交给算子启动程序执行，完成算子代码的加载和启动，并测试其耗时。DeepGen 以 PyTorch 作为基准对代码结果进行正确性验证；如果验证通过，则根据方案耗时决定是否将其纳入最优算子缓存。
\begin{figure}[htbp]
\centering
\includegraphics[width=0.9\textwidth]{3_31}
\caption{DeepGen 自动调优流程}\label{3_31}
\vspace{\baselineskip}
\end{figure}


% =============================================================
%                  4.2 调优空间参数的设置与定义
% =============================================================
\subsection{调优空间参数的设置与定义}

在生成算子代码的过程中，主要调优参数会决定中间表示的变换路径，并影响最终代码的性能表现。这些参数主要包括线程块和线程的数据分块尺寸、软件流水线阶段设置、局部规约轴切分策略，以及线程块、线程束和线程的布局结构。上述因素共同构成了高维且具有离散特征的非凸联合优化空间。测试结果说明，即使处于相同硬件平台，仅对少量参数进行变更也可能造成性能大幅改变。在各类硬件架构上，参数配置的最优设置也会呈现不同表现。采用固定参数的传统优化方法，已经难以适应不同环境下 Transformer 模型对高性能的需求。为此，本文形成了一套同 IR 优化流程紧密结合的调优参数方法，用于为后续自动调优模块提供结构化、可扩展的参数空间定义。

\begin{table}[htbp]
\caption{调优参数列表}\label{tab:table5}
\vspace{0.5em}\centering\wuhao
\renewcommand{\arraystretch}{1.35}
\begin{tabular}{cc}
\toprule[1.5pt]
参数类型 & 描述 \\ \midrule[1pt]
block\_size & 线程块所操作的数据分块尺寸，包括行、列和规约轴三个方向 \\ 
warp\_size & 线程束所需要计算的数据分块尺寸，包括行、列两个方向 \\ 
thread\_size & 线程所需要计算的数据分块尺寸，包括行、列两个方向 \\ 
scatter\_block\_size & 离散化后线程块需要计算的数据分块尺寸，包括行、列两个方向 \\ 
scatter\_warp\_size & 离散化后线程束需要计算的数据分块尺寸，包括行、列两个方向 \\ 
group\_size\_m & 线程块索引重映射分组的行方向线程块个数 \\ 
pipeline\_stage\_num & 软件流水线阶数，包括线程和线程块两个级别，1表示不启用优化 \\ 
local\_split\_num & 局部规约轴切分次数，1表示不启用优化 \\ 
unroll & 循环展开次数 \\ \bottomrule[1.5pt]
\end{tabular}
\vspace{\baselineskip}
\end{table}

如表 \ref{tab:table5} 所示为本研究设计的调优参数集合。
其中，block\_size、warp\_size 与 thread\_size 共同定义了内核启动的线程网格与线程块结构，以及线程块内线程束与线程的布局关系。
\begin{figure}[htbp]
\centering
\includegraphics[width=0.9\textwidth]{3_32}
\caption{注意力机制调优配置参数}\label{3_32}
\vspace{\baselineskip}
\end{figure}
这三个参数共同构建了 IR 的基础结构，并直接决定了算子在 GPU 内存层次中的数据流动与计算效率。
scatter\_block\_size 与 scatter\_warp\_size 为 AccessScatter Pass 提供配置依据，通过调整线程块与线程束的内存访问粒度，有效规避共享内存的 Bank 冲突。
group\_size\_m 参数用于 BlockSwizzle Pass，通过选择合适的分组大小提升数据局部性，从而提高缓存命中率与执行效率。
pipeline\_stage\_num 用于控制软件流水线的深度设置，MultiBuffer Pass 依据流水线阶数扩展循环中具有依赖关系的缓冲区，并由 SoftwarePipeline Pass 完成相应的 IR 变换。
local\_split\_num 参数控制 LocalSplitK Pass 对局部规约轴进行切分的次数，通过增加线程数量以提升并行度。
unroll 参数则用于配置循环展开次数，通过减少循环控制开销与增强指令级并行以优化性能。
\begin{figure}[htbp]
\centering
\includegraphics[width=0.9\textwidth]{3_33}
\caption{矩阵乘法调优配置参数}\label{3_33}
\vspace{\baselineskip}
\end{figure}

图 \ref{3_32} 和图 \ref{3_33} 表明，该策略提供一个可配置的 JSON 接口，以此对调优参数的搜索空间进行定义。用户利用配置文件能够灵活指定各参数候选取值的范围，自动调优系统依据此信息生成完整的参数组合空间供搜索使用。为了提高参数组织的结构化程度，本文对调优参数进行了划分，形成二级参数方法。在数据分块类参数中，二级参数对应各个计算方面的分块尺寸；在软件流水线类参数中，二级参数对应各层级循环的流水线深度设置。

图 \ref{3_32} 展示了注意力算子的调优配置示例。由于该算子内部包含两个存在数据依赖关系的矩阵乘操作，两者共享相同的线程块尺寸约束，因此 block\_size 参数仅设置一种二级配置。对于第二个矩阵乘操作，其第二输入矩阵需要单独设置离散化访存参数，因此 scatter\_block\_size 与 scatter\_warp\_size 均包含两组二级参数配置，其中第二组配置的列方向参数留空，表示继承前一矩阵乘操作的列方向设置。
图 \ref{3_33} 展示了通用矩阵乘法的配置示例。相较于注意力算子，其额外包含 group\_size\_m 参数，这是因为注意力算子受其计算逻辑的约束，无法应用 BlockSwizzle Pass 优化。


% =============================================================
%                  4.3 搜索空间剪枝策略
% =============================================================
\subsection{搜索空间剪枝策略}
在自动调优过程中，算子的性能一般会由多个调优参数来决定。本文以线程块级数据分块尺寸、线程级计算粒度、规约轴切分策略、软件流水线阶段数、线程块映射方式等参数为例进行分析。各参数之间存在比较复杂的耦合关系，不同参数的搭配组合对应着完全不同的执行策略。从理论上看，对所有可能的参数组合进行穷举搜索，可以得到最优的实现方案。不过，实际应用时，调优参数一般会形成一个高维离散空间，该空间的大小会随着参数数量的增多而以指数级方式增长。如果直接对全部候选方案进行编译和测试，则会产生很高的时间开销和计算资源消耗，并影响自动调优效率。

\begin{figure}[htbp]
\centering
\includegraphics[width=0.95\textwidth]{3_34}
\caption{调优空间剪枝过程}\label{3_34}
\vspace{\baselineskip}
\end{figure}

为了提高调优效率，本文在生成调优空间之后引入一种分层剪枝方式，该方式在保证算子功能正确和硬件可以执行的基础上，尽可能减少实际需要进行编译和测试的候选方案数量，整体流程如图 \ref{3_34} 所示。剪枝策略按照作用方式的区别可以分成两类，包括依靠硬件资源约束和架构特性的必要剪枝、依靠经验规则和性能启发的可选剪枝。前者用于保证方案可行，后者用于进一步压缩搜索空间。

初始调优空间生成完成以后，所有参数组合虽然都符合基本的生成条件，但是其中仍然存在大量可以在运行环境中执行但性能潜力偏低的配置。本文首先运用可选剪枝策略，将明显不具备优势的方案提前过滤。例如在数据复制过程中，每个线程为了保证访存效率，一般会负责一定规模的数据搬运工作。如果拷贝的数据量很少且只使用单线程，线程调度和同步开销就会增大，进而影响执行效率。矩阵乘算子里，分块形状会影响数据复用效果。当计算分块接近正方形时，数据局部性更好；当搜索空间较大时，系统会优先保留该分块组合。

在完成可选剪枝后，调优策略接着进行必要剪枝，让剩余方案可以在目标硬件上正确运行。本部分主要根据 GPU 的资源限制和执行模型展开。共享内存和寄存器资源都有上限，当某组参数让资源需求超过硬件限制时，需要进行剔除。线程块的组织方式也需要满足线程束的基本要求，否则会影响硬件利用率。如果共享内存访问方式不合理，则会导致 Bank 冲突，进而降低访存带宽利用率。

将前面提到的两种剪枝策略结合使用后，原始调优空间可以得到有效压缩，调优过程主要分布在潜在的高性能区域中。当搜索空间较小时，系统可以只启用必要剪枝操作；而在搜索空间较大时，系统可以同时启用可选剪枝操作，从而降低搜索复杂度并提高调优效率。

% 第五节
\section{本章小结}
本章详细阐述了面向 Transformer 模型关键算子的代码生成与 DeepGen 的核心设计与实现细节。针对现有编译器在块级逻辑表达和跨平台性能迁移方面的局限性，本章从中间表示构建、编译优化变换以及自动化调优三个维度提出了一套系统化的解决方案。

首先，本文提出了一种面向异构硬件层次结构的多级中间表示建模方法。本文设计了名为 DeepGen 的方言，把算子的计算逻辑在线程块级和线程级两个层次上进行解耦，并设计了降级信息推断算法，解决了将粗粒度块级表示转换为细粒度线程级表示的索引映射问题，从而能够准确描述 GPU 计算调度和访存路径。

其次，在分析矩阵乘法和注意力机制的算术及访存特性后，本章设计并实现了一组针对性的高性能优化 Pass 集合。具体内容包括：通过软件流水线隐藏全局内存和共享内存访问延迟；通过局部规约轴切分增加长规约轴情况下的线程级并行度；通过线程块访存重排提高 Cache 复用率；通过离散化访存消除共享内存的 Bank 冲突。多个 Pass 协同工作，缓解了算子执行过程中的内存与计算瓶颈。本章还针对硬件平台差异与移植策略进行了分析，指出不同硬件在存储层次、线程调度和执行粒度上的差异如何影响优化 Pass 的实际收益，并阐述了将统一 IR 与平台相关参数约束结合的必要性，以确保生成代码能够适配不同硬件平台。

最后，为了解决多维优化参数组合爆炸问题，本文形成了利用空间剪枝的自动调优框架。本文通过定义结构化的调优参数方法，并结合硬件资源约束下的必要剪枝和专家经验指导的可选剪枝策略，剔除无效、低效的参数组合，从而提高调优效率。该策略实现了编译、测试和监控的异步协同，在缩小搜索空间的同时保证调优过程的高效性。
